\section{Síntesis y ampliación}

La línea central de esta clase cabe en una sola frase: el crecimiento de las capacidades de la IA no es la curva de un solo modelo, sino una transformación de pila completa impulsada a la vez por context, compute, capital y culture. Context determina con qué tareas y retroalimentación reales entra en contacto el sistema; compute determina la escala de experimentación que puede sostenerse; capital determina con qué rapidez se construye la infraestructura; culture e institutions determinan los objetivos, las responsabilidades y la asignación de recursos. Si se descuida cualquiera de las capas, la expansión de otra puede convertirse en desperdicio o en riesgo.

\begin{figure}[H]
\centering
\includegraphics[width=0.92\textwidth]{images/01-03-45-food-for-thought.jpg}
\caption{Las dos preguntas que el curso deja a los estudiantes: cómo lograr una compute transition pacífica y de qué manera puede participar cada persona.}
\figsource{Video oficial de Stanford Online, 01:03:45, `images/01-03-45-food-for-thought.jpg`.}
\end{figure}

\begin{importantbox}{Lectura de la figura: el cierre no es una pregunta de predicción para espectadores}
La primera pregunta exige desglosar la «transición pacífica» en accesibilidad, interoperabilidad, concentración del mercado, energía, cadenas de suministro transfronterizas, seguridad e interés público; la segunda pide a los estudiantes que propongan acciones que ellos mismos puedan verificar, por ejemplo definir estándares de recursos, construir workloads portátiles, publicar la metodología de un benchmark, investigar la planificación y la liquidación, participar en organismos de estandarización o explicar las consecuencias de una política. El valor del proyecto del curso no está solo en mostrar un modelo, sino en que los estudiantes entren en la conversación real sobre infraestructura.
\end{importantbox}

\begin{knowledgebox}{Nota de clase: considérate un active participant}
Al final, el ponente enfatiza que los estudiantes no necesitan llegar a ser CEO, directores de investigación o reguladores para tener derecho a participar. Describir con claridad los problemas de una interfaz, publicar experimentos reproducibles, documentar migraciones fallidas, proponer borradores de estándares y pedir evidencia a los ponentes de la industria son acciones que pueden cambiar la discusión. La meta de la preparedness no es predecir mejor quién ganará, sino ser capaz de formular mejores preguntas y construir sistemas verificables en un periodo de incertidumbre.
\end{knowledgebox}

\subsection{Ocho conclusiones transferibles}

\begin{enumerate}
  \item Primero traza las dependencias de toda la pila y después discute los avances en una sola capa; los modelos, los productos y la capacidad de cómputo están insertos en un sistema de recursos más amplio.
  \item Trata la scaling law como un modelo empírico: precisa el intervalo observado, los supuestos de extrapolación y las condiciones en que deja de cumplirse.
  \item Desglosa el context en tareas, herramientas, entorno, permisos y retroalimentación; no lo entiendas solo como un prompt largo.
  \item Al auditar RL, revisa al mismo tiempo el environment, la reward, el verifier y los límites de responsabilidad.
  \item Al leer casos de la industria, sigue el punto de entrada, el modelo, el entorno de ejecución, el retorno de datos y el control de la distribución.
  \item En las gráficas macroeconómicas, distingue las señales tecnológicas, de adopción, de capital y de precios; no uses un tipo de evidencia para sostener una conclusión de otro tipo.
  \item Al usar analogías históricas, transfiere los mecanismos, no calendarios definidos; revisa en especial la vida útil de los activos y las condiciones de estandarización.
  \item Para que el compute se convierta en una infraestructura ampliamente accesible, deben madurar juntos los estándares técnicos y las instituciones que los hacen cumplir.
\end{enumerate}

\subsection{Lecturas adicionales}

Para seguir aprendiendo, se sugieren tres rutas reproducibles. Primera: consulta el sitio oficial del curso Stanford CS153 y los videos oficiales de Stanford Online, y compara cómo los invitados posteriores modifican el marco de cuatro restricciones de esta clase. Segunda: elige un coding agent o una tarea de entrenamiento y traza su context flow, su entorno de ejecución, sus señales de verificación y sus autorizaciones de datos. Tercera: compara la job specification, el SLA y los criterios de facturación de dos mercados de nube o de GPU, y comprueba en la práctica si un workload puede migrarse, en lugar de comparar solo los precios de lista.

\begin{knowledgebox}{Recordatorio de las notas: cómo usan esta clase las sesiones posteriores}
Con cada invitado, registra al menos cuatro cosas: la capa del sistema en la que trabaja, el recurso que considera más escaso, el tipo de evidencia que presenta y los problemas que espera que resuelva el mercado o las instituciones. Al terminar las diez semanas, reúne estas respuestas en una misma tabla; los puntos de conflicto suelen enseñar más que los consensos.
\end{knowledgebox}

\subsection{Tarea práctica: auditar una ruta real de Compute}

Para convertir el marco macro de esta clase en una capacidad de sistemas, puedes elegir un servicio de IA que uses realmente y rastrear una solicitud desde su origen hasta los recursos de cómputo. El objetivo de la tarea no es adivinar toda la arquitectura interna del proveedor, sino dejar claro qué hechos son observables, cuáles provienen de documentación pública y cuáles siguen siendo inferencias. El entregable final debe permitir que otra persona revise las interfaces de recursos, el retorno del context y los límites de responsabilidad.

\begin{enumerate}
  \item Traza la ruta de la solicitud desde el cliente, el servicio de la aplicación y la API del modelo hasta el clúster de inferencia, e indica en qué capa se aplican la identidad, los datos y los permisos de las herramientas.
  \item Elige una tarea repetible y registra la latencia, los tipos de falla, la calidad de la salida y el comportamiento de los reintentos; explica si la reward o la condición de éxito puede verificarse automáticamente.
  \item Compara las especificaciones, los precios y los pasos de migración de dos canales de suministro, y revisa si ``GPU-hour'', el token price o el nombre de una API ocultan que algo es insustituible.
  \item Propón una mejora mínima de estándar, por ejemplo una job spec portátil, códigos de falla unificados, divulgación del rendimiento o comprobantes de eliminación de datos, y explica qué organismo impulsaría su adopción.
\end{enumerate}

\begin{importantbox}{Criterios de aceptación}
Una tarea aceptable debe incluir enlaces a las fuentes, pasos experimentales reproducibles, al menos un caso de falla, columnas separadas para hechos e inferencias, y una explicación de su impacto en la seguridad y el interés público. Si la conclusión solo puede formularse como «tal empresa es más fuerte» o «se necesita más capacidad de cómputo», el análisis todavía no ha llegado a las capas de interfaces, mecanismos e instituciones que exige Frontier Systems.
\end{importantbox}

\end{document}
